%% math1-qf-common-0001 — 二次関数の最大・最小(共通テスト風・独自問題)
%% Manabi Forge contributors による独自制作の非公式教材。
%% 過去問題・公式サンプル問題の複製や言い換えではない。
%% コンテンツライセンス: CC BY 4.0(リポジトリルートの LICENSE-CONTENT)
\documentclass[common-test]{manabi}

\ManabiId{math1-qf-common-0001}
\ManabiVersion{0.1.0}
\ManabiTitle{二次関数の最大・最小(演習)}
\ManabiCourse{数学I}
\ManabiUnit{二次関数}
\ManabiDifficulty{標準}
\ManabiMinutes{15}

\begin{document}
\ManabiMakeTitle

\begin{mnbproblem}[(配点 20)]{1}
太郎さんは、長さ $20\,\mathrm{m}$ のフェンスをちょうど使い切って、
長方形の花壇を作ることにした。花壇の 4 つの辺すべてをフェンスで囲むとき、
隣り合う 2 辺の長さを $x\,\mathrm{m}$ と $(10 - x)\,\mathrm{m}$ と表せる。
花壇の面積を $y\,\mathrm{m}^2$ とする。

\begin{mnbparts}
  \item $y$ を $x$ の式で表すと
    \[
      y = -x^2 + \mnbans{アイ}\, x
    \]
    である。また、$x$ のとり得る値の範囲は $0 < x < \mnbans{ウエ}$ である。

  \item $y$ は $x = \mnbans{オ}$ のとき最大値 $\mnbans{カキ}$ をとる。

  \item 通路を確保するため、$x$ の範囲を $2 \le x \le 4$ に制限する。
    このとき、$y$ は $x = \mnbans{ク}$ のとき最大値 $\mnbans{ケコ}$ をとる。

  \item (3) の範囲 $2 \le x \le 4$ において、$y$ が最小となる $x$ の値として
    正しいものを、次の \mnbcircled{1} ~ \mnbcircled{4} のうちから一つ選べ。
    解答欄: \mnbans{サ}

    \begin{mnbchoices}
      \item $x = 2$
      \item $x = 3$
      \item $x = 4$
      \item 最小となる $x$ の値はない
    \end{mnbchoices}
\end{mnbparts}
\end{mnbproblem}

\ManabiSolutionPart

\begin{mnbparts}
  \item 長方形の周の長さが $20\,\mathrm{m}$ なので、隣り合う 2 辺の和は
    $10\,\mathrm{m}$ である。よって 2 辺は $x$ と $10 - x$ であり、
    \[
      y = x(10 - x) = -x^2 + 10x.
    \]
    辺の長さは正であるから $x > 0$ かつ $10 - x > 0$、すなわち
    $0 < x < 10$。\quad(\mnbans{アイ} $= 10$、\mnbans{ウエ} $= 10$)

  \item 平方完成すると $y = -(x - 5)^2 + 25$。グラフは上に凸で頂点は
    $(5,\ 25)$、頂点は定義域 $0 < x < 10$ に含まれる。よって
    $x = 5$ のとき最大値 $25$。\quad(\mnbans{オ} $= 5$、\mnbans{カキ} $= 25$)

  \item 頂点の $x$ 座標 $5$ は範囲 $2 \le x \le 4$ に含まれない。
    この範囲では $y$ は増加するから、右端 $x = 4$ で最大となり
    $y = -(4)^2 + 10 \cdot 4 = 24$。\quad(\mnbans{ク} $= 4$、\mnbans{ケコ} $= 24$)

  \item (3) と同様に、$2 \le x \le 4$ で $y$ は増加するから、最小は左端
    $x = 2$ でとり $y = -(2)^2 + 10 \cdot 2 = 16$。正解は \mnbcircled{1}。
    \quad(\mnbans{サ} $= 1$)

    \textbf{誤答の傾向}: \mnbcircled{3} は最大と最小の位置の取り違え、
    \mnbcircled{2} は「中央で最小」という誤解、\mnbcircled{4} は閉区間で
    端点が最小値を与えることの見落としによる。
\end{mnbparts}

\end{document}
